\section{Part Nine: A Diagnostic Ladder for a Single Wrong Output}

The prior parts decompose a \emph{tool} or a \emph{task}. A complementary decomposition applies to a single \emph{output} --- the moment something a system produced is simply wrong, and the people involved need a precise, shared vocabulary for what kind of wrong it is.

\textbf{The five axes, ordered from most mechanical to most human:} structure, meaning, context, groundedness, environment.

\textbf{Structure: is the output well-formed?} A structurally correct output satisfies the format's formal constraints --- JSON that parses, code that imports, a query the database accepts. The most mechanical axis, and the one modern systems handle best; the fix lives in the decoding stack (schema-validated sampling, constrained generation). The diagnostic trap runs opposite the failure rate: because structural correctness is so easy to verify, practitioners over-rely on it and stop checking once it passes. A syntactically perfect object can still assert something false or refer to something that doesn't exist.

\textbf{Meaning: is the output internally coherent?} A structurally sound output can still contradict itself --- a paragraph whose first sentence contradicts its third, a function whose declared return type doesn't match what it returns. These failures are detectable through careful reading alone, without checking against the outside world.

The underlying mechanism, stated precisely: a model trained to predict the next token given everything before it optimizes for the \emph{local} plausibility of each individual token. A token can be highly plausible at every position and the sequence can still, as a whole, arrive at a global contradiction, because nothing in a single generation pass architecturally checks the output against itself as a complete artifact. Reasoning-oriented training and inference-time techniques (sampling multiple times and checking consistency, self-critique, a separate verification pass) narrow this gap without closing it, at the cost of additional compute.

\textbf{Context: is the output appropriate for the situation, given the information it was actually given?}

\textbf{Groundedness: do the output's factual claims connect to verifiable reality, or faithfully reflect the sources it was given?}

\textbf{A structural correction to how v1 distinguished context from environment.} v1 defined both axes principally as failures caused by missing situational information in the prompt, which made them descriptively almost identical and undermined the ladder's usefulness as a diagnostic tool --- if two axes point at the same underlying mechanism, walking "through" both doesn't actually add diagnostic resolution. The cleaner, non-overlapping distinction: \textbf{a context failure is a failure to correctly use or weigh information that was actually available to the system} --- the audience, tone, or purpose was stated or inferable from what was provided, but the model didn't apply it correctly (dense jargon delivered despite a stated non-specialist audience; a five-paragraph essay delivered despite a request for a quick yes-or-no). \textbf{An environment failure is a failure caused by information that was never made available to the system at all} --- the user knew something relevant (an internal convention, a specific meaning of a domain term, a piece of context from outside the conversation) and simply never wrote it into the prompt, so no amount of correct use of the \emph{given} information could have produced the right answer.

This distinction matters practically: a context failure is fixed by better use of what's already there --- often a harness-level or prompting-technique fix. An environment failure can only be fixed by supplying the missing information in the first place; no amount of better instruction-following inside the system helps if the fact simply isn't present anywhere in its context.

\textbf{Returning to groundedness, in full:} groundedness failures are especially hard to catch because no internal reading of the output reveals them --- a fabricated citation has the same structural form as a real one. The check has to come from outside the text.

\textbf{A correction to how v1 characterized the relationship between model confidence and reliability --- this is the single most important correction in this document, since v1 stated it in a way that was both too strong and internally inconsistent with an adjacent sentence.} v1 claimed that a model's expressed confidence is unrelated to its actual reliability, and that confidently-stated claims are "exactly as likely" to be wrong as hedged ones. \textbf{That's factually incorrect, and it contradicted v1's own acknowledgment two sentences earlier that calibration research shows model confidence correlates with accuracy above chance.} The corrected, internally consistent version: \textbf{expressed confidence is often meaningfully miscalibrated and must never be treated as proof of correctness on its own --- but where a specific model's confidence signal has been independently validated for a specific task, it may be used as one input to risk-based triage, not as the sole basis for deciding what to check.} The practical difficulty, and the reason the stricter reading still mostly holds in practice: that validation work is task- and model-specific, is rarely done, and the signal even where validated is not strong enough to substitute for external verification of anything that actually matters. So the working rule that remains defensible: \textbf{treat any consequential factual claim as an unverified hypothesis by default, and only rely on a confidence signal for triage where that signal's calibration has actually been measured for the model, task, and deployment in question} --- which in most real deployments, absent that measurement work, collapses back to "check everything that matters" in practice.

\textbf{When groundedness failures cluster in a specific domain or deployment context, the root cause is frequently distributional} --- the training distribution (whether general-purpose or narrowly fine-tuned) doesn't match the deployment distribution. Retrieval and tool use reduce individual errors but don't address this underlying mismatch; the remedy is targeted evaluation against representative deployment-distribution examples, not incremental per-call fixes.

\textbf{Environment: was the question itself fully specified?} The hardest axis to see, because it's invisible from the model's side by construction. The user knows things the model has no access to; if those things aren't written into the prompt, the model cannot use them. The diagnostic signature: the user reads the answer and feels frustrated even though the answer, on its own terms, is fine --- it just doesn't address the real situation. The fix is to write down what the model needed to see, not to ask again with the same information.

\textbf{A correction to how the ladder-walking procedure was stated.} v1 instructed walking the ladder bottom-up and stopping at "the first axis" that contains a failure. \textbf{That's unsafe as a general procedure.} Outputs commonly have more than one simultaneous failure mode, and a structural pass says nothing about whether a normative, factual, or contextual problem also exists in the same output. The corrected procedure: \textbf{use the bottom-up order as an efficient, low-cost sequence for checking --- cheaper checks first, since they're more likely to explain an issue at lower diagnostic cost --- but continue through all five axes and record every applicable failure mode before closing the diagnosis, rather than stopping at the first hit.} Lower-axis failures remain, on average, easier to diagnose and cheaper to fix, which is exactly why checking them first is efficient --- but "checked first" is not the same claim as "sufficient once found."

\textbf{Chains and intermediate outputs.} For any task chaining multiple steps, apply the ladder at each intermediate output, not only the final result. \textbf{A correction to an internal inconsistency in v1's original statement of this point:} v1 claimed a chain where every step individually "passes the five-axis walk" can still propagate an early false premise forward undetected. As stated, that's self-contradictory --- if an early output genuinely passed a sufficiently rigorous \emph{groundedness} check, the false premise wouldn't have survived that check to propagate. The corrected, coherent version: \textbf{local structural and coherence checks at an intermediate step can all pass while that step's groundedness was never actually checked} --- a function can be well-formed and internally consistent while calling a method that doesn't exist, and unless groundedness is specifically re-checked at that checkpoint (not just structure and meaning), the false premise passes forward as an accepted input to the next step. The remedy is checkpointing with \textbf{specifically groundedness-inclusive} intermediate verification, not just structural or coherence checks at each stage.

\textbf{Beyond the five axes, four further failure classes escape the single-output frame entirely, each requiring a different kind of remediation:}

\emph{Normative} --- an output can be correct by all five axes and still wrong by a value standard the user holds. Requires governance and value-alignment work, not a better single-output check.

\emph{Adversarial} --- an output deliberately induced by a malicious prompt, poisoned retrieval source, or manipulated tool, with indirect prompt injection (instructions embedded in retrieved content rather than the user's own prompt) as the practically important variant. The underlying model has no reliable mechanism to distinguish an instruction from data. \textbf{A correction to v1's proposed remedy: v1 presented "input sanitization, output schema enforcement, and privilege separation" as if this were a complete defense.} It isn't --- these are components of a defense-in-depth posture that also needs least-privilege access, action confirmation before consequential steps, content-provenance tracking, and ongoing monitoring, not a fixed short list that, once implemented, closes the risk.

\emph{Drift} --- an output correct yesterday, wrong today because the world changed; a monitoring and governance problem, not a single-output problem.

\emph{Complacency} --- the ladder assumes an actively engaged reviewer. Automation complacency and automation bias are well-supported human-factors phenomena --- reduced detection of failures and missed anomalies under highly reliable automation is a documented pattern. \textbf{A correction to how precisely v1 stated the consequence:} v1 claimed the effective error rate "rises steadily" while the model's own error rate stays constant, phrased as a near-universal monotonic trajectory. The underlying research supports the general pattern (vigilance decays, and systems typically record whether review happened rather than how carefully) more than it supports a specific claim about a steadily rising trajectory in every AI workflow --- the corrected version keeps the mechanism (decaying vigilance under a run of correct-seeming outputs, tracked imperfectly by systems that record participation but not attentiveness) without asserting a specific universal shape to the resulting curve.
